\section{Discussion}
\label{sec:discussion}

\paragraph{What this work shows.}
SR-Core provides a hard working-set guarantee: exactly $k$ blocks are active per token,
regardless of bank size or recursion depth.
Entropy-based router consolidation exposes a controllable axis within this fixed working
set: increasing $\lambda$ moves the router toward more concentrated, stable routing paths,
reducing simulated bytes-in-motion and unique routing vocabulary.
At $\lambda = 0.003$ we observe no measured tradeoff: cache locality and held-out code
generalization both improve relative to the unregularized baseline, and the improvement
replicates across two independent seeds.
At $\lambda = 0.005$ the cache benefit is larger and the quality effect is seed-dependent.
Dense models remain the quality upper bound throughout.

\paragraph{What this work does not show.}
We do not demonstrate wall-clock inference speedup.
The dispatch overhead of sparse block selection imposes a CPU-side penalty (approximately
$2.8\times$ slower than compute-matched dense in our CPU benchmark), so the
bytes-in-motion reduction does not directly translate to latency improvement without
hardware co-design.
We also do not demonstrate quality parity with dense models, scale the bank beyond
$n = 64$ blocks, or evaluate on downstream tasks.
The bytes-in-motion estimates are from simulation under idealized LRU caching; real
transfer costs depend on block size, DRAM bandwidth, PCIe characteristics, and prefetch
scheduling, none of which are accounted for here.

\paragraph{Scope of the entropy objective.}
Entropy minimization sharpens the pre-selection distribution at $r = 1$; it does not
explicitly maximize diversity or cross-token overlap.
This means it is not a general solution to the reuse/novelty tension identified in
Section~\ref{sec:tension}: it moves the router toward the reuse end of the axis in a
controlled way, which happens to benefit cache efficiency.
Whether a different objective that explicitly targets diversity at deep recursion steps
could yield complementary gains remains an open question.

\paragraph{Future directions.}
The most immediate validation gap is real-latency measurement: an end-to-end
offloading prototype (RAM\,$\to$\,VRAM) would connect the bytes-in-motion simulation to
observable tokens/second.
Scaling the bank size beyond $n = 64$ is also necessary to test whether the working-set
invariant holds at the regime (large $n$, very small active fraction) where the
transfer-reduction ratio is most compelling.
On the regularization side, targeted entropy objectives (e.g., penalizing only when entropy
exceeds a threshold) and their interaction with load-balancing terms remain unexplored at
the trained-model scale.
